% =========================================================================
% Related work fragment for "Frozen Basis Time Series Forecasting"
% (camera-ready rewrite of NeurIPS submission 30981).
%
% Governing documents: paper_rewrite_plan.md §7 subplan and Part II §19
% (corrections entries 3 and 10); response_K8yb.md Weakness (i) and
% response_ovPT.md Question 1 (Pramanick comparison and caveats);
% notation_and_taxonomy_spec.md §7 rule 3.
%
% All citations verified against v2_OxAlpha_rewrite/references.bib.
% Oreshkin et al. is cited as 2020 throughout (plan §12).
% =========================================================================

\section{Related Work}
\label{sec:related}

\subsection{What a frozen basis can and cannot carry}
\label{sec:related:approximation}

Approximation theory separates two regimes that the submitted version
conflated. A fixed orthonormal basis performs \emph{linear}
approximation: the quality of a rank-truncated expansion depends on how
much of the signal's energy the retained rows span, and among
orthonormal bases the Karhunen--Lo\`eve transform is optimal for that
compaction, with the discrete cosine transform close behind on realistic
windows \citep{mallat1989wavelet}. Wavelets earn their reputation in the
other regime, \emph{nonlinear} approximation, where selection of the
best coefficients adapts to each signal \citep{daubechies1992ten}. A
basis fixed at build time never adapts per series, so it lives in the
first regime regardless of its family, which is why our localization
prior inverted: the wavelet advantage is a theorem about adaptive
selection, and at build time Karhunen--Lo\`eve wins by construction.
This reframing also answers, honestly and without new experiments, where
the alternatives sit: short-time Fourier, wavelet-packet,
stationary-wavelet, and empirical-mode constructions differ in what
generates the basis rows, but captured energy governs any of them once
rows are fixed, so we did not run them as separate arms.

The identity result connects this lineage to random bases. At complete
square rank any two orthonormal bases span the same space and differ by
a rotation that the learned projection absorbs, which we verified at
residual $5.6 \times 10^{-15}$ and measured across wavelet, cosine and
random orthonormal draws; random draws therefore serve as controls
rather than as baselines to beat. Compression through constrained
bottlenecks has its own long history in deep learning
\citep{hinton2006reducing, vincent2008denoising}; our contribution there
is the measurement that in N-BEATS the trunk, not the structured head,
is what the parameters were doing.

\subsection{Wavelets inside deep forecasters}
\label{sec:related:waveletml}

Prior work wires fixed transforms into deep forecasters along three
routes: as input preprocessing feeding conventional networks
\citep{deznabi2023multiwave}, as transforms embedded in attention or
Transformer components \citep{sasal2022wtransformers, zhou2025wavets},
and as basis expansions paired with N-BEATS itself
\citep{pramanick2024fusion}. The submitted version differentiated itself
from the last route by assertion alone. It is now differentiated by
measurement, and the distinction came out narrower than asserted.

\subsection{The measured comparison with DWT preprocessing}
\label{sec:related:pramanick}

The design of \citet{pramanick2024fusion} applies a discrete wavelet
transform to the input window and trains a separate generic stack per
sub-band; ours puts a frozen basis inside each block with no
preprocessing. We implemented their design from the published
description and compared the two under one protocol: 420 runs, three
matched budgets, six M4 periods, ten matched seeds, arms paired by
actual parameter count. The in-block frozen basis wins all ten
budget-matched comparisons on the median (sign test $p = 0.0020$) and
nine of ten on the mean, by margins of $+0.083$ to $+1.308$. At the
larger rung on Quarterly, Weekly and Daily their arm carries 23--24\%
more parameters and still loses, and against the generic baseline their
design never wins on any period. Three caveats belong beside those
numbers. The single cell where the mean favours their design turns on
one collapsed seed of our own arm, and the median still favours ours
there. Two of our four registered predictions failed, including the one
predicting parity between the designs. And this is our reimplementation
without their code, on data their paper did not use, so it is good-faith
work rather than an independent replication. A fourth regime was
discarded before analysis: rectifying every stack is structurally unfair
to a sub-band design, because a detail band is a zero-mean oscillation
whose corrections must cancel, so the activated regime tests the
rectifier rather than the placement question.

The narrowed claim is therefore twofold. In-block frozen bases beat
sub-band preprocessing at matched budget, and neither design's accuracy
comes from the basis, since our own controls show the basis inert at
complete rank. The originality claim shrinks accordingly, and this paper
states it at that size.

\subsection{The N-BEATS lineage}
\label{sec:related:nbeats}

N-BEATS \citep{oreshkin2020nbeats} established basis expansion with
doubly residual stacks as a single-model architecture competitive with
ensembles on M4 \citep{makridakis2020m4}, and NHiTS
\citep{challu2023nhits} kept the block interface while replacing the
schedule with multi-rate pooling and hierarchical interpolation. Our
results read back into both. Against N-BEATS, the interpretable
configuration carries roughly fifty times more trunk capacity than its
single-model accuracy requires, a measurement about the architecture
that no member of the lineage had made, with the caveat that the
published ensemble results are unswept. Against NHiTS, the transfer
study shows the block interface carries effects across the backbone
change while the tuning does not, which supports the interface design
and cautions against importing hyperparameters with it.


